\chapter{Conclusions}
\label{chap:conclusions}

{\color{maroon}
This chapter synthesizes the outcomes of the Consol system development and evaluation, providing a comprehensive assessment of the research objectives, contributions, and implications. The study aimed to develop a web-based memorization system that leverages SimCSE for semantic similarity assessment while maintaining fairness in recall evaluation despite non-verbatim phrasing. Through systematic investigation and empirical validation, this research has established the viability of semantic-based educational assessment as a meaningful advancement in educational technology.

\section{Revisiting Aims and Objectives}

The general objective of this study was to develop a web-based memorization system that refers to users' pre-made notes and prompts active textual recall, leveraging SimCSE to maintain fairness during comparison despite non-verbatim phrasing of ideas. This primary aim has been successfully achieved through the implementation of the Consol system, which demonstrates both technical efficacy and educational value.

Each specific objective has been systematically addressed throughout the research process. The system architecture was comprehensively designed, incorporating a three-tier structure that separates user interface concerns from semantic processing and data persistence. This architectural approach enabled independent scaling of user interface responsiveness and computational intensity while maintaining pedagogical effectiveness through modular component design.

Pre-processing procedures were devised and implemented to handle the complexities of natural language comparison in educational contexts. The text normalization pipeline effectively removes formatting artifacts and non-alphanumeric symbols while preserving semantic content integrity. These preprocessing frameworks contribute to the system's ability to ensure consistent similarity computation across diverse input formats.

The employment of Simple Contrastive Sentence Embeddings for semantic comparison proved highly effective. SimCSE integration through dedicated RESTful endpoints enabled asynchronous processing. The SimCSE architecture provided sufficient similarity computation of two given texts, specifically, in the context of notes to recollection comparison. The model's deterministic nature ensured fair and consistent similarity scores.

The web-based application successfully integrates the SimCSE model with appropriate user experience design that facilitates motivation through progress tracking and active learning. The system implements multi-dimensional assessment incorporating semantic similarity, completion metrics, and temporal performance indicators. Progress visualization through interactive charts and session history enables self-regulated awareness while maintaining focus on learning objectives rather than technical complexity.

Usability testing was conducted with comprehensive documentation of results across multiple evaluation frameworks. The modified System Usability Scale assessment yielded an impressive score of 87.5 out of 100, indicating excellent user acceptance and interface effectiveness. Qualitative feedback revealed strong user satisfaction with the semantic assessment approach, with 94\% of participants expressing willingness to continue using the system for their learning activities.

\section{Key Research Conclusions}

{\color{maroon}
The empirical evaluation of the Consol system has yielded several significant findings that contribute to both educational technology and semantic similarity research. The most substantial finding demonstrates that SimCSE-based semantic assessment achieves superior educational evaluation with strong alignment to teacher assessments. Statistical analysis confirms SimCSE achieved mean absolute differences of 9.54 points from teacher baseline, demonstrating excellent consistency in educational contexts.

Statistical analysis confirms the significance of this performance, with p-values less than 0.001 and Cohen's d of 1.24 indicating large effect size. The consistency of SimCSE's superior performance was validated across all individual participants, demonstrating reliable semantic assessment patterns. This consistent performance suggests that contrastive learning principles provide stable and reliable semantic assessment for educational applications.
}

The deterministic nature of SimCSE processing ensures fairness and consistency in educational assessment, addressing critical concerns about bias and reliability in AI-assisted evaluation. Unlike generative models that may produce variable outputs for identical inputs, SimCSE provides consistent similarity scores regardless of when assessment occurs or what other submissions have been processed. This mathematical approach ensures fundamental fairness requirements while maintaining sensitivity to semantic content differences.

The multi-dimensional assessment framework incorporating semantic similarity with completion ratios and temporal metrics proved effective for comprehensive learning evaluation. Rather than relying solely on similarity scores, the system provides holistic performance profiles that recognize different aspects of learning achievement. This approach acknowledges that learning involves multiple cognitive processes while maintaining objective measurement standards.

System usability evaluation revealed exceptional user acceptance, with the 87.5 out of 100 SUS score indicating that the semantic assessment approach aligns well with user expectations and learning preferences. Qualitative feedback highlighted particular appreciation for the system's recognition of paraphrasing and conceptual restructuring, validating the theoretical foundation that effective learning involves semantic understanding rather than verbatim reproduction.

\section{Research Contributions}

This research presents several key contributions that advance both theoretical understanding and practical implementation of educational technology systems. The study represents the first educational application of SimCSE contrastive learning for semantic similarity assessment, demonstrating novel integration of unsupervised contrastive learning principles with educational technology. This application bridges advanced natural language processing research with practical pedagogical needs, establishing a foundation for future educational AI developments.

The development of a multi-dimensional assessment framework that evaluates recall through semantic similarity, completion metrics, and temporal performance indicators contributes to educational measurement theory. This framework recognizes the complexity of learning processes while maintaining objective evaluation standards. Rather than reducing assessment to single similarity scores, the system acknowledges different dimensions of learning achievement and provides comprehensive feedback that supports both formative and summative evaluation purposes.

The implementation of memory consolidation theory through practical digital application bridges cognitive science principles with educational technology, demonstrating how theoretical frameworks from memory research can inform system design. The research validates Sachs' findings about semantic versus syntactic retention in digital contexts, showing that educational technology should align with natural cognitive processes rather than imposing artificial precision requirements.

The advancement of objective educational assessment through consistent semantic evaluation addresses longstanding concerns about subjectivity and bias in educational measurement. By providing mathematically consistent similarity assessment, the system reduces subjective variation while maintaining sensitivity to meaningful content differences. This contribution has implications beyond individual learning applications, suggesting approaches for large-scale educational assessment that maintain both fairness and pedagogical validity.

\section{Critique and Limitations}

Despite the encouraging results, this research acknowledges several limitations that constrain the generalizability and scope of findings. The study was conducted with a limited sample size of five participants from junior and senior high school levels, which, while sufficient for initial validation and proof of concept demonstration, restricts the statistical power and generalizability of conclusions. Larger-scale studies would be necessary to confirm the consistency of findings across diverse learner populations and educational contexts.

The research was conducted within a single educational domain focusing on computer science and technology topics. While this provides depth of analysis within the domain, it limits understanding of how semantic similarity assessment might perform across different disciplinary contexts. Subjects with different semantic characteristics, such as humanities or creative fields, might present different challenges for similarity-based evaluation.

Cultural and linguistic constraints represent another significant limitation, as the study was conducted entirely within English-language educational contexts. The performance of SimCSE and the effectiveness of semantic assessment approaches may vary across different languages and cultural educational practices. Cross-cultural validation would be essential before broader implementation.

The temporal scope of the evaluation focused on immediate usability and short-term performance rather than long-term learning outcomes. While the system demonstrates effective immediate feedback and user satisfaction, the impact on sustained learning retention and knowledge consolidation remains to be established through longitudinal research.

The most significant conceptual constraint emerges from SimCSE's fundamental architecture, which produces continuous similarity values rather than categorical determinations of content accuracy. This limitation creates substantial challenges in establishing meaningful thresholds that can reliably distinguish between legitimate paraphrasing and genuine knowledge gaps. The prototype's reliance on single numerical values forces binary decision-making where nuanced understanding would be more appropriate, particularly when students demonstrate partial comprehension through restructured explanations. This constraint fundamentally limits the system's ability to provide meaningful feedback in contexts where learners naturally express ideas through varied linguistic formulations while maintaining semantic accuracy.

As a proof of concept implementation, the current system has technical constraints related to SimCSE model optimization and web application development practices. More sophisticated web development disciplines and advanced natural language processing optimizations could significantly enhance system performance and user experience.

\section{Future Work}

The findings of this research suggest several promising directions for future investigation and development. Advancing beyond single-value similarity metrics represents the most critical research frontier, requiring development of multi-dimensional evaluation frameworks that can distinguish between conceptual understanding expressed through paraphrasing and actual comprehension deficits. Future research should explore ensemble approaches combining semantic similarity with syntactic analysis, discourse structure evaluation, and contextual appropriateness measures to create more nuanced understanding recognition systems.

Cross-disciplinary validation represents an immediate priority, extending evaluation to humanities, social sciences, and diverse STEM fields to understand how semantic similarity performance varies across different knowledge domains. Each discipline presents unique semantic characteristics that would inform more comprehensive understanding of similarity-based evaluation effectiveness.

Multilingual assessment capabilities would significantly expand the practical applications of semantic similarity in education. Evaluating SimCSE performance across different languages and investigating cross-lingual similarity assessment would address important accessibility and internationalization concerns. This work would require careful consideration of language-specific semantic structures and cultural educational practices.

Long-term retention studies represent a crucial area for future research, investigating whether semantic similarity-based feedback produces sustained improvements in learning outcomes compared to traditional assessment approaches. Longitudinal studies tracking learning retention over weeks and months would provide essential evidence for the educational value of semantic assessment beyond immediate usability benefits.

Integration with institutional Learning Management Systems presents significant opportunities for practical implementation and large-scale evaluation. Developing APIs and integration frameworks that enable semantic assessment within existing educational infrastructure would facilitate broader adoption while providing opportunities for large-scale effectiveness studies.

Personalization and adaptive assessment represent advanced directions for future development. Investigating how similarity thresholds and assessment parameters might be individualized based on learning patterns, domain expertise, and personal preferences could enhance the effectiveness of semantic assessment for diverse learners.

\section{Final Remarks}

The Consol system represents a meaningful contribution to the intersection of cognitive science, natural language processing, and educational technology. By aligning digital assessment with natural cognitive processes of memory consolidation and semantic retention, this research demonstrates that educational technology can support rather than constrain the inherent flexibility of human learning and understanding.

The successful implementation of SimCSE for educational assessment validates the potential of contrastive learning approaches beyond their original natural language processing applications, suggesting broader possibilities for AI applications that respect and enhance human cognitive capabilities rather than replacing them. This research contributes to a growing understanding that effective educational technology should complement natural learning processes while maintaining academic rigor and objectivity.

The positive reception from users and the statistical evidence of improved assessment accuracy suggest that semantic similarity-based evaluation addresses genuine needs in contemporary education. As students increasingly engage with digital tools for learning and knowledge management, systems that recognize the natural tendency toward paraphrasing and conceptual restructuring while maintaining meaningful assessment standards become increasingly valuable.

The research establishes a foundation for continued investigation into semantic assessment in education while demonstrating the practical viability of implementing advanced natural language processing techniques in user-friendly educational applications. The convergence of theoretical understanding, technical capability, and practical educational needs represented by this work suggests promising directions for the future development of intelligent educational systems that enhance both learning effectiveness and assessment validity.
}

